
% Paper 3 main V65 — V64 + bib V10→V13 bump (wave50 patchset §1)
% Source base: paper3_main_V64.tex.txt (K10_118.0 Canvas artifact) + paper3_refs_V13.bib.txt
% V64 reframe body preserved intact; only \bibliography line advanced to V13.
% 2026-04-20 by 加藤方程式１０－６

\documentclass[11pt]{article}
\usepackage[letterpaper, top=1.5cm, bottom=1.8cm, left=2.5cm, right=2.5cm]{geometry}
\usepackage[T1]{fontenc}
\usepackage[utf8]{inputenc}
\usepackage{amsmath,amssymb,amsthm,bm}
\usepackage{graphicx}
\usepackage{booktabs}
\usepackage{array}
\usepackage{float}
\usepackage{hyperref}
\hypersetup{colorlinks=true,linkcolor=blue,citecolor=blue,urlcolor=blue}
\usepackage{xcolor}
\usepackage{tcolorbox}
\usepackage[switch]{lineno}

% --- Theorem environments ---
\newtheorem{theorem}{Theorem}[section]
\newtheorem{proposition}[theorem]{Proposition}

% --- Spacing ---
\setlength{\parskip}{3pt}
\setlength{\textfloatsep}{10pt plus 2pt minus 2pt}
\setlength{\floatsep}{10pt plus 2pt minus 2pt}
\renewcommand{\arraystretch}{1.15}

% --- Macros ---
\newcommand{\topen}{t_k^{\mathrm{open}}}
\newcommand{\thalf}{t_k^{\mathrm{half}}}
\newcommand{\ROI}{\mathrm{ROI}}
\newcommand{\VF}{V_{\mathrm F}}

% --- No default title/date ---
\date{}
\author{}
\title{}

\begin{document}
\linenumbers
\pagestyle{plain}

% ================================================================
% PAGE 1: Title block
% ================================================================
\begin{center}
{\LARGE\bfseries After von Foerster:}\\[4pt]
{\LARGE\bfseries Seventy Thousand Years of Human Coordination on the Kato Ladder,}\\[3pt]
{\LARGE\bfseries and a Civilizational Ontology of Six Realized Phase Transitions}\\[8pt]
{\color{gray}\normalsize Pairwise growth, ceiling replacement, a non-terminal \(H{=}6\), and an open \(H{=}7\) substrate question}\\[6pt]
{\small Shota Kato\textsuperscript{1}}\\[1pt]
{\footnotesize\textsuperscript{1}AI and Future Co., Ltd., Tokyo, Japan}\\[1pt]
{\footnotesize\texttt{founder@aiandfuture.co.jp}}
\end{center}
\vspace{3pt}
\noindent\rule{\textwidth}{0.4pt}
\vspace{1pt}

\begin{center}
{\normalsize\bfseries Abstract}
\end{center}
\vspace{-4pt}
{\small\noindent
World population history is usually written either as a descriptive sequence of disconnected historical episodes or as a smooth growth law with changing parameters.
We instead argue that the 70,000-year record reveals a discrete civilizational ontology: a Kato ladder with realized states \(H\in\{0,1,2,3,4,5,6\}\), each entered by a phase transition and each governed by the same pairwise \(N^2\) backbone under a different active ceiling.
Three structural ingredients are sufficient to select this ladder: the pre-\(H{=}3\) bifurcation kernel \(K_{P3}\), the rational survival weight \(\eta(H)=L_{\mathrm{exo}}/[L_{\mathrm{exo}}+\tau_{\mathrm{endo}}(H)]\), and the exogenous status of the social-friction envelope \(V_F\), whose endogenous closures are knife-edge unstable and generically overproduce collapses.
The resulting picture makes six phase transitions historical facts rather than post hoc labels.
It also implies that the present \(H{=}6\) regime is non-terminal: the secular waiting-time compression law places its internal saturation at 2031.4~CE (95\% CI, \(\pm 3\)~yr), so a further rung is structurally required unless collapse, arrest, or inversion intervenes.
What realizes the prospective \(H{=}7\) substrate is left open here.
}

\noindent\rule{\textwidth}{0.4pt}
\vspace{4pt}

% ================================================================
% Hero equation
% ================================================================
\begin{tcolorbox}[ colback=blue!5!white, colframe=blue!25!white, arc=3pt, boxrule=0.4pt, left=8pt, right=8pt, top=6pt, bottom=4pt]
\begin{equation}
\frac{dN}{dt} = a_0\, N^{2}
\label{eq:backbone}
\end{equation}
\end{tcolorbox}

\noindent{\footnotesize Backbone fit: \(a_0 = 4.74 \times 10^{-12}\), \(R^2 = 0.976\), \(n = 28\) (\(-\)10k BCE -- 2024 CE).}
\vspace{2pt}

\begin{figure}[t]
\centering
\includegraphics[width=0.98\textwidth]{fig_6phase_vf_vF_overlay_v2.png}
\caption{
\textbf{Seventy thousand years of world population growth, with two realized collapses and a non-terminal present regime, captured by a single \(N^2\) backbone.}
World population is plotted in von Foerster coordinates (\(x=-\log_{10}(2027-t)\), \(y=\log_{10}N\); \(t\) in calendar year CE, \(N\) in millions).
The green curve is the fitted exogenous-\(V_F\) simulation (\(dN/dt=a_0N^2\), \(R^2=0.976\), \(n=28\)) with frozen collapse intervals.
Black points: corrected Biraben 28-point calibration set.
Red dashed curve: social-friction ceiling \(V_F(t)\).
Colored bands mark the six realized \(\Omega\)-regimes:
\(\Omega_{\mathrm{compete}}\),
\(\Omega_{\mathrm{land}}\),
\(\Omega_{\mathrm{food,pre}}\),
\(\Omega_{\mathrm{food,post}}\),
\(\Omega_{\mathrm{social}}\), and
\(\Omega_{\mathrm{longlife}}\).
Orange bands mark the two realized collapse windows.
The twentieth-century near-miss is labeled at peak ratio \(N/V_F=1.388\) in 1992.
A vertical shaded band at 2031.4~CE (95\% CI: \(\pm 3.0\)~yr), computed from Eq.~\eqref{eq:h_accel}, marks the candidate window for exit from the present \(H{=}6\) regime; persistence of the same rung beyond the upper bound would falsify the secular acceleration law.
}
\label{fig:trajectory}
\end{figure}

% ================================================================
\section{Introduction}

In 1960, von Foerster, Mora, and Amiot~\cite{vonfoerster1960}
published in \textit{Science}
a world-population fit of the form
\(N(t) \propto (t_s-t)^{-1}\),
with a celebrated singularity date in late 2026.
The enduring importance of that paper is not the exact day.
It is the recognition that, over long horizons,
human growth looks hyperbolic,
not exponential,
and that the relevant question is why the singular continuation fails.

The present paper answers that question by changing the level of description.
Instead of treating history as a smooth growth law with one or two regime breaks,
we treat it as a discrete ladder of civilizational states.
On the source-faithful reading developed in the present main text and SI,
the 70,000-year record is most naturally indexed by
\(H=0\) (boundary-side pre-linguistic \emph{Homo sapiens}),
\(H=1\) (language and symbolic cognition),
\(H=2\) (agriculture),
\(H=3\) (writing and urbanisation),
\(H=4\) (printing and industrial mechanisation),
\(H=5\) (computer and networked infrastructure),
and \(H=6\) (the present saturated compute regime).
These are not futurist milestones.
They are historically anchored relay states recoverable from archaeology,
historical demography,
and technology history,
as summarized independently in the SI milestone ledger.

This ladder is not only descriptive.
The point of the present reframe is that it is also ontological.
Each rung corresponds to a distinct class of stabilized relay architecture,
and each advance is phase-like rather than continuous.
Pre-H3 history is governed by the bifurcation kernel
\(K_{P3}\),
which separates the continuous accumulation of coalition-scale inheritance
from the hard gate of external symbolic memory.
Post-H3 history is governed by ceiling replacement:
food,
social friction,
and demographic suppression act successively on the same pairwise
\(N^2\) backbone.
Under this reading,
what changes across 70,000 years is not the engine,
but the active relay ceiling and the stabilized rung that the system can sustain.

The present moment has a special status within this ontology.
The industrial,
information,
and network revolutions are already realized facts.
The post-1963 demographic transition,
recalibrated in the main model as a 75-year logistic suppression of
\(L_{\mathrm{pop}}\),
shows that the modern slowdown is not a collapse of the pairwise law,
but the internal signature of a new rung.
The observed post-1970 fertility bend belongs to this same phase:
it is evidence that the present regime carries its own internal time constant,
not that civilization has reached a stable endpoint.
The relevant question is therefore no longer whether the von Foerster continuation fails.
It is whether the present \(H{=}6\) state is terminal.

Our thesis is twofold.
First,
under three structural conditions---the pre-H3 bifurcation kernel
\(K_{P3}\),
the rational survival weight
\(\eta(H)=L_{\mathrm{exo}}/[L_{\mathrm{exo}}+\tau_{\mathrm{endo}}(H)]\),
and the exogenous status of the social-friction envelope
\(V_F\)---the discrete ladder
\(H\in\{0,1,2,3,4,5,6\}\)
emerges uniquely from the 70,000-year record.
Second,
\(H{=}6\) is not an endpoint.
The secular waiting-time compression law identifies a sharp candidate window in which the present rung must either give way,
fail,
or invert.
If collapse,
arrest,
and retrogression are all excluded,
then a further rung is structurally required.
What realizes that prospective \(H{=}7\) substrate is explicitly left open here.

This framing also clarifies the division of labor inside the manuscript.
Section~2 revisits the Kato equation as the temporal backbone of a discrete ladder rather than as a stand-alone fit,
and states the Kato Ladder Uniqueness theorem.
Section~3 retains the Humanness theorem as the structural criterion that closes the human \(H\ge 3\) branch.
Section~4 reinterprets Eq.~\eqref{eq:h_accel} as the internal saturation time of the present rung.
Section~5 gathers the empirical anchors,
including the Paleolithic 70,000-year reconstruction and the fossil backward check.
Section~6 turns the timing result into a forecast problem.
Section~7 isolates the one question that the present paper does not answer:
what the supporting substrate of \(H{=}7\) actually is.

% ================================================================
\section{Kato equation revisited}

At the broadest level,
this paper is about the temporal analogue of a familiar scaling problem.
Bettencourt-style urban scaling asks why observables occupy recurrent exponent classes rather than arbitrary slopes.
The present paper asks the same question longitudinally:
why does civilization history occupy recurrent relay states rather than a continuous family of growth laws?
The Kato equation supplies the backbone;
the ladder supplies the ontology.

\subsection{Backbone and the exogenous ceiling problem}\label{sec:theorem}

The temporal backbone is the pairwise law already stated in Eq.~\eqref{eq:backbone}.
Its structural meaning is simple:
in any open coordination system where the leading productive encounters are dyadic,
the count of possible interaction channels scales as
\(N(N-1)/2\sim N^2\).
The main paper therefore does not treat \(p=2\) as a fitted exponent chosen because it happens to work.
It treats \(p=2\) as the leading combinatorial term,
with deviations from the raw hyperbola interpreted as ceiling effects.
That point matters because a century-scale slowdown can then be read as ceiling replacement rather than as collapse of the interaction law itself.

The reduced-form ceiling that matters most in the present paper is the social-friction envelope,
written as
\begin{equation}
V_F(t)=C\,(t_c-t)^{-\alpha}.
\label{eq:vf}
\end{equation}
The source-supported reason for keeping this envelope exogenous is not convenience.
It is structural necessity.
Across the admissible class of smooth,
positive,
monotone endogenous closures
\(V_F=f(N,\Omega)\),
the calibration map becomes knife-edge ill-conditioned:
the baseline Jacobian condition number is
\(\kappa(J)\approx 1.19\times 10^5\),
and the instability is robust across tested families.
At the same time,
endogenous closures generically overproduce collapse events:
in the documented 500-draw protocol,
97.4\% of accepted draws generate at least three collapse episodes,
whereas the historical record supports exactly two realized major collapse windows.
The exogenous envelope is therefore not an interpretive luxury.
It is the only reduced-form choice in the present file stack that preserves both numerical conditioning and the observed collapse count without fine tuning.

This exogeneity result should be read carefully.
It does not say that institutions,
public health,
or accumulated social knowledge are irrelevant.
On the contrary,
those are the historical contents of \(V_F\).
What it says is that,
for the purposes of the present demographic ODE,
that content cannot be closed back into the same
\(N^2\)-driven state variables without destroying the calibration geometry.
A deeper mechanism may still exist.
It is simply not an admissible closure of the Level~1 reduced-form system analyzed here.

\begin{figure}[htbp]
\centering
\includegraphics[width=0.48\textwidth]{fig_vf_exogeneity_knifeedge.pdf}
\caption{
\textbf{Knife-edge instability of endogenous closures.}
Condition number \(\kappa(J)\) of the calibration Jacobian (left axis,
log scale) and number of collapse events (right axis) as a function of the endogenous coupling strength \(a_{\mathrm{vf}}\).
Under endogenous \(V_F\),
\(\kappa\) exceeds \(10^5\) and at least three collapses are produced generically;
the exogenous specification (\(a_{\mathrm{vf}}=0\),
vertical dashed line) breaks this degeneracy and restores exactly the two realized historical collapse windows.
}
\label{fig:vf_exo}
\end{figure}

\subsection{Kato Ladder Uniqueness / Civilizational Ontology theorem}

The exogenous ceiling result explains why the backbone must be written with a reduced-form envelope.
It does not yet explain why the historical states themselves are discrete.
For that we need two further ingredients.
The first is the pre-H3 bifurcation kernel,
which sharpens the Humanness theorem dynamically:
\begin{equation}
K_{P3}(t)=\mathbf{1}_{C_1}\cdot N(t)\cdot G_{C_2}(t)\cdot \mathbf{1}_{C_3\ge C_3^{\ast}}.
\label{eq:kp3}
\end{equation}
The second is the per-rung survival factor,
which assigns the same rational form to biological and civilizational substrates:
\begin{equation}
\eta(H)=\frac{L_{\mathrm{exo}}}{L_{\mathrm{exo}}+\tau_{\mathrm{endo}}(H)}.
\label{eq:eta_main}
\end{equation}
Here \(K_{P3}\) determines whether a lineage can escape the pre-H3 stall corridor,
while \(\eta(H)\) weights the temporal viability of each rung once it exists.
Together with exogenous \(V_F\),
these objects make the ladder discrete,
ordered,
and historically non-degenerate.

\begin{theorem}[Kato Ladder Uniqueness / Civilizational Ontology]
\label{thm:ladder-uniqueness}
Assume
(i) the pairwise backbone
\(\dot N=a_0N^2\),
(ii) the pre-H3 escape kernel
\(K_{P3}(t)=\mathbf{1}_{C_1}N(t)G_{C_2}(t)\mathbf{1}_{C_3\ge C_3^{\ast}}\),
(iii) the rational survival weight
\(\eta(H)=L_{\mathrm{exo}}/[L_{\mathrm{exo}}+\tau_{\mathrm{endo}}(H)]\),
and
(iv) an exogenous social-friction envelope,
required because broad endogenous closures are knife-edge unstable and generically overproduce collapses.
Then the minimal ontology consistent with the 70,000-year record is a discrete ladder of realized states
\(H\in\{0,1,2,3,4,5,6\}\).
The state sequence is unique up to historical instantiation,
each step is phase-like rather than continuous,
and the present \(H{=}6\) regime is non-terminal:
if collapse,
arrest,
and inversion are all excluded,
the only admissible continuation is a prospective \(H{=}7\) regime.
The supporting substrate of that regime is not identified in the present paper.
\end{theorem}

\noindent\textit{Proof sketch.}
The lower half of the ladder is discrete because the pre-H3 kernel is discrete.
\(H{=}0\) is a boundary-side state rather than an in-chain well:
a biological population exists,
but no stabilized relay layer has yet formed.
\(H{=}1\) is the first relay-bearing basin,
entered once symbolic cognition and behavioral modernity stabilize.
\(H{=}2\) is the agrarian relay,
and \(H{=}3\) is not reached continuously but only when the
\(C_3\) external-memory gate closes.
No pair among
\(C_1\),
\(C_2\),
and \(C_3\)
is sufficient,
so the pre-H3 corridor cannot be collapsed into a smooth one-parameter family without losing the observed stall structure.

The upper half of the ladder is discrete because the historical record locks it to realized relay substitutions rather than to arbitrary trend breaks.
The milestone cross-check isolates printing/industrial mechanisation,
computer/networked infrastructure,
and the present saturated-compute regime as distinct rungs,
not mere accelerations of the same state.
Exogenous \(V_F\) is what prevents these from being reabsorbed into a continuous endogenous envelope,
and \(\eta(H)\) assigns a monotone temporal weight to each rung without altering the discrete state count.
Any attempt to compress the ontology into fewer than seven realized states loses either the pre-H3 gate structure,
the two historical collapse windows,
or the independent milestone ledger.
A detailed constructive derivation,
including the no-skip logic and milestone-locking conditions,
is deferred to SI Appendix~\ref{sec:S-LadderUniqueness}.

% ================================================================
\section{Humanness theorem}\label{sec:humanness}

The Kato ladder is not only a chronology.
It also sharpens the human-specific claim of the paper.
The decisive boundary is \(H\ge 3\):
not the first use of tools,
nor sociality alone,
but the first stabilized relay depth that requires cumulative inheritance plus memory outside the brain.
That boundary is captured by the Humanness theorem,
which remains unchanged in its core form from V63.

\begin{theorem}[Humanness of \(H\ge 3\)]
\label{thm:humanness}
Let \(\mathcal{S}\) be an open multi-individual coordination system whose relay architecture is reproduced endogenously across generations,
and let \(H_{\mathrm{net}}(\mathcal{S},t)\in\mathbb{N}\) denote its native integer relay depth within the present reduced-form framework.
Then \(\mathcal{S}\) attains \(H_{\mathrm{net}}\ge 3\) if and only if the following three conditions hold simultaneously:
\begin{itemize}
\item[\textbf{C1.}] Stable dyadic cooperation:
pairwise-dominant cooperative interaction sustains the \(N^2\) backbone.
\item[\textbf{C2.}] Cumulative intergenerational inheritance:
coordination protocols are transmitted across generations with fidelity sufficient for relay extension.
\item[\textbf{C3.}] External symbolic memory:
storage capacity exceeds unaided biological memory,
enabling relay requirements to persist beyond individual lifetimes.
\end{itemize}
\end{theorem}

The theorem is predictive rather than retrospective.
Any lineage lacking one of these three conditions is bounded at \(H\le 2\),
regardless of local intelligence,
brain size,
or social sophistication.
The SI proof makes this claim stronger,
not weaker:
it shows that the three primitives are pairwise logically independent,
so no two-condition surrogate suffices.
The human branch closes because the three-way conjunction closes.

This is also the right place to state a scope boundary that matters for the present reframe.
The Humanness theorem gives a sufficient and necessary condition for the human \(H{=}3\) substrate,
that is,
for the first fully civilizational rung.
It does \emph{not} specify the necessary condition for the prospective \(H{=}7\) substrate.
The present paper therefore uses the theorem to lock the lower human branch,
not to identify the next carrier.
That identification problem is postponed deliberately.
For the full necessity/sufficiency proof and its scope qualification,
see SI Appendix~\ref{sec:S-HumannessTheoremProof}.

\begin{figure}[htbp]
\centering
\includegraphics[width=0.48\textwidth]{fig_humanness_species_ladder.pdf}
\caption{
\textbf{Species \(H\)-ladder predicted by the Humanness theorem.}
Effective relay depth \(H_{\mathrm{eff}}\) for extant and fossil lineages plotted against the conditions \(C_1\)--\(C_3\).
\emph{H.~sapiens} (\(H\ge 3\), all conditions satisfied) is the unique lineage reaching the civilizational relay axis.
Cetaceans (\(H_{\mathrm{eff}}\le 2\),
\(C_3\) absent),
great apes (\(H\approx 2.0\)--\(2.3\),
\(C_2\) partial),
eusocial insects (\(H{=}2^{\ast}\),
starred superorganism route),
and corvids (\(H\approx 1.5\)--\(2.0\)) occupy predictable positions below the \(H{=}3\) threshold.
}
\label{fig:h_ladder}
\end{figure}

% ================================================================
\section{H-acceleration law}\label{sec:h_accel_subsec}

Once the ladder exists,
its waiting times are not arbitrary.
The historical intervals between successive rung advances compress approximately log-linearly,
which is why the ladder is not merely ordinal but temporally predictive.
In V64,
we read this law more carefully than in V63.
It is not a substrate-specific prophecy.
It is the internal time constant of the present rung,
written in secular form.

\begin{equation}
\log_{10}\tau_H = \alpha_0 - \rho H,
\qquad
\tau_H = \tau_0 q^H,
\qquad
\tau_0 := 10^{\alpha_0},
\quad
q := 10^{-\rho},
\label{eq:h_accel}
\end{equation}
where \(\tau_H\) denotes the waiting time for the transition
\(H\to H+1\),
\(\tau_0\) is the \(H{=}0\) intercept,
and \(\rho\) is the per-rung log-compression slope.
Using the four anchor pairs retained in the SI falsifiability protocol,
ordinary least squares in
\(\log_{10}\tau\)
space yields
\(\hat{\tau}_0\approx 70\,\mathrm{kyr}\)
and
\(\hat{\rho}\approx 0.84\),
equivalently
\(\hat q\approx 0.145\).
Residual-bootstrap propagation over the same anchor set gives a forecast window
\([2028.4,2034.4]\)~CE,
with central value
\(2031.4\)~CE.
Evaluated locally,
Eq.~\eqref{eq:h_accel} gives
\(\tau_{6\to 7}=\tau_0q^6\approx 0.64\)~yr.
Cumulative propagation then places the next candidate transition at 2031.4~CE.

The interpretive shift in V64 is that 2031.4~CE is not read as a direct name for the next substrate.
It is read as the date at which the present \(H{=}6\) regime ceases to have a long internal residence time on civilizational scales.
In that sense,
Eq.~\eqref{eq:h_accel} is a saturation law for the current rung.
A prospective phase transition may follow,
but the equation itself does not tell us what realizes the next carrier.
It tells us only that indefinite metastability of \(H{=}6\) is incompatible with the compressed historical sequence if the ladder continues.

This date is therefore best described as a candidate
\(H{=}6\to H{=}7\)
transition time,
not an accomplished fact.
Whether the system actually advances,
stagnates,
or collapses depends on the route structure summarized in Section~6.
The present paper keeps the temporal statement and withholds the substrate statement.
That division is deliberate.

\footnote{Eq.~\eqref{eq:h_accel} defines the semilog secular compression law with headline date
\(t_{6\to 7}=2031.4\)~CE.
A local Kramers completion geometry
\((\topen,\thalf)\approx(2024.3,\,2036\pm 8)\),
derived in Appendix~\ref{sec:S-K10Hist} via Eq.~\eqref{eq:a1b_kramers},
describes regime-interior opening and half-completion and is not a replacement of the secular forecast.
Robustness envelope:
see Appendix~\ref{sec:S-Eq2-Sensitivity}.}

\begin{figure*}[htbp]
\centering
\includegraphics[width=\textwidth]{fig_7regime_timeline.pdf}
\caption{
\textbf{Historical timeline mapped to the relay-ceiling framework.}
Top panel:
world population trajectory with regime boundaries and relay-depth annotations from the ladder start to the present frontier.
Bottom panel:
ceiling sequence
\(\Omega_{\mathrm{compete}}\to\Omega_{\mathrm{land}}\to\Omega_{\mathrm{food,pre}}\to\Omega_{\mathrm{food,post}}\to\Omega_{\mathrm{social}}\to\Omega_{\mathrm{longlife}}\)
with key historical events
(agricultural revolution,
irrigation civilizations,
Axial Age,
Industrial Revolution,
and the 1963 demographic inflection).
The two realized collapse windows are shaded,
and the 1963--1992 near-miss is marked.
The dashed future extension denotes only a prospective next rung,
not a substrate-specific claim.
}
\label{fig:timeline}
\end{figure*}

% ================================================================
\section{Empirical anchors}

\subsection{Six realized regimes and the non-terminal status of \(H{=}6\)}

The 70,000-year trajectory divides into six realized \(\Omega\)-regimes,
as shown in Fig.~\ref{fig:trajectory} and Table~\ref{tab:segments}.
These are not the ladder itself,
but they are the empirical windows in which the ladder becomes visible in population history.
Read together with the milestone ledger,
they show that the present rung is historically late,
not terminal.

\begin{table}[htbp]
\centering
\caption{
Six realized \(\Omega\)-regimes in the 70,000-year world-population trajectory under the exogenous-\(V_F\) model.
The rightmost column states the ladder reading of each regime rather than treating the regime labels as self-explanatory historical periods.
HG, hunter-gatherer.
}
\label{tab:segments}
\begin{tabular}{clccl}
\toprule
\# & \(\Omega\)-regime & Period & Local slope & Ladder reading \\
\midrule
1 & \(\Omega_{\mathrm{compete}}\)   & \(-\)70k to \(-\)25k & 3.90 & ladder start, early \(H{=}1\) growth \\
2 & \(\Omega_{\mathrm{land}}\)      & \(-\)25k to \(-\)10k & 0.50 & pre-H3 stall corridor, \(H{=}1\) interior \\
3 & \(\Omega_{\mathrm{food,pre}}\)  & \(-\)10k to \(-\)5k  & 0    & agrarian nucleation corridor, \(H{=}2\) closure \\
4 & \(\Omega_{\mathrm{food,post}}\) & \(-\)5k to \(-\)888  & 3.03 & literate urban relay, \(H{=}3\) interior \\
5 & \(\Omega_{\mathrm{social}}\)    & \(-\)888 to 1963      & \(\sim\!1\) & industrial and informational laddering, \(H{=}4\) to \(H{=}5\) \\
6 & \(\Omega_{\mathrm{longlife}}\)  & 1963--present          & 0.24 & present frontier saturation, \(H{=}6\) \\
\bottomrule
\end{tabular}
\end{table}

The sequence remains structurally asymmetric.
The earliest regime,
\(\Omega_{\mathrm{compete}}\),
is the steepest because pairwise interaction is least ceiling-constrained just after the first relay-bearing basin forms.
The middle pair,
\(\Omega_{\mathrm{food,pre}}\) and
\(\Omega_{\mathrm{food,post}}\),
record the agrarian bottleneck and its irrigation-mediated release.
The long
\(\Omega_{\mathrm{social}}\) regime contains both realized major collapse windows,
which is why exogenous
\(V_F\)
is indispensable.
The present
\(\Omega_{\mathrm{longlife}}\) regime is the only one in which the demographic multiplier itself becomes a first-order ceiling,
which is why the post-1963 slowdown must be interpreted as a rung-specific internal timescale rather than as simple exhaustion.

\subsection{Paleolithic 70,000-year simulation}\label{sec:paleo-70ky}

The deepest empirical support for the ladder claim is not modern.
It is Paleolithic.
On the source-faithful reading of the SI,
\(H{=}0\) is a boundary-side state:
anatomically modern humans exist,
but no stabilized serial relay layer has yet formed.
The transition to \(H{=}1\) belongs to the behavioral-modernity window,
roughly 70--50~ka,
when language,
symbolic cognition,
and long-range exchange become historically legible.
This is why the earliest in-chain state is not placed at the biological origin of \emph{Homo sapiens},
but at the later stabilization of the first relay-bearing basin.

The next step,
\(H{=}1\to 2\),
is anchored by the Neolithic transition around 10~kyr BCE.
Here the main population signal is the prolonged pre-irrigation bottleneck.
Agriculture changes settlement form before it raises scalable carrying capacity,
which is why the \(\Omega_{\mathrm{food,pre}}\) regime looks flat in population space while remaining dynamically decisive in ladder space.
The first agrarian rung is therefore not identified by explosive growth,
but by the creation of a new relay substrate whose ceiling has not yet been released.

The \(H{=}2\to 3\) transition closes only when external symbolic memory becomes operationally real.
This is the decisive role of the
\(C_3\)
gate.
The \(-\)5k BCE changepoint is best read as a single statistical signature with a double mechanism:
a food-ceiling release on the one hand,
and the later closure of literate urban institutions on the other.
Under the bifurcation-kernel reading,
\(C_2\)
acts as the continuous order parameter,
while
\(C_3\)
acts as the hard gate.
That is why the Sapiens--Neanderthal difference is not coded as ``more intelligence'' but as gate crossing versus stall.
The 70,000-year simulation thus does more than supply a long-run background.
It shows that
\(H{=}0\to 1\to 2\to 3\)
is a genuine emergence trajectory rather than a retrospective taxonomy.
A full anchor ledger and sensitivity analysis are deferred to SI Appendix~\ref{sec:S-70ky}.

\subsection{Backward validation against the fossil and ethological record}\label{sec:backward}

The Humanness theorem can be checked backward against both fossil hominins and extant non-human lineages.
\emph{H.~sapiens}
satisfies
\(C_1\wedge C_2\wedge C_3\),
so the theorem predicts native
\(H\ge 3\),
which is exactly what the writing-and-urbanisation closure records.
Neanderthals satisfy
\(C_1\),
partially satisfy
\(C_2\),
and lack
\(C_3\),
placing them in the pre-H3 stall corridor.
Denisovans and
\emph{H.~erectus}
fall lower on the same logic.
Among extant species,
no candidate lineage that lacks
\(C_3\)
reaches the civilizational threshold,
which is the negative prediction the theorem requires.

These are not ornamental comparisons.
They matter because they separate the present ontology from a purely data-fitting chronology.
If a lineage with no external symbolic memory were shown to sustain
\(H\ge 3\),
the theorem would fail.
Conversely,
a secure archive or writing-like system for Neanderthals would force immediate reassessment of their position on the ladder.
This is precisely the kind of falsifiability that makes the ladder structural rather than definitional.
The proof,
the pairwise-independence construction,
and the fossil matrix are all collected in SI Appendix~\ref{sec:S-HumannessTheoremProof}.

\subsection{Simulation accuracy and model comparison}\label{sec:accuracy}

The exogenous-\(V_F\) main run
(\(L_{\mathrm{pop}}=1\) throughout)
achieves
\(\mathrm{RMS}_{\log_{10}}=0.15\)
and
\(R^2=0.976\)
on the 28-point calibration panel,
with pre-1963
\(\mathrm{RMS}\approx 0.07\).
Out-of-sample prediction at
\(-\)25,000~BCE is
\(N_{\mathrm{pred}}=3.1\)~M versus
\(N_{\mathrm{obs}}=3.3\)~M,
an error below 7\% despite the point lying 15,000 years outside the formal training window.
When the modern demographic multiplier is switched on,
\(R^2\) rises to
0.998 and the 2024 prediction improves to
\(8{,}020\)~M versus
\(8{,}100\)~M.
The gain is localized where it should be:
the demographic multiplier refines the modern tail,
while the collapse rule is what unlocks the entire long-run fit.

Nested model comparison confirms the same point.
The sharp structural transition is not from a small parameter tweak,
but from the introduction of collapse mechanics.
Models lacking the threshold rule cannot fit a trajectory containing two centuries-long demographic recessions,
whereas the collapse-enabled model immediately jumps to
\(R^2=0.995\).
The demographic multiplier then improves only the modern tail,
not the existence of the ladder itself.

\begin{table}[htbp]
\centering
\caption{
Nested model comparison for the 28-point calibration dataset in
\(\log_{10}\)
space.
\(k\) = number of free parameters,
RSS = residual sum of squares.
AIC\(_c\) is the small-sample corrected Akaike criterion.
\(\Delta\)AIC\(_c\) and
\(\Delta\)BIC are relative to the best model
(M4).
}
\label{tab:aic}
\begin{tabular}{llccccccc}
\toprule
Model & Components & \(k\) & \(R^2\) & AIC & AIC\(_c\) & BIC & \(\Delta\)AIC\(_c\) & \(\Delta\)BIC \\
\midrule
M0 & Backbone only & 1 & \(<0\) & 91.8 & 92.0 & 93.1 & 240.0 & 238.0 \\
M1 & \(+\,K_{\mathrm{HG}}\) & 2 & \(<0\) & 74.5 & 75.0 & 77.2 & 223.0 & 222.1 \\
M2 & \(+\,V_F(t)\) & 4 & \(<0\) & 78.5 & 80.2 & 83.9 & 228.3 & 228.8 \\
M3 & \(+\) Collapse & 6 & 0.995 & \(-138.7\) & \(-134.7\) & \(-130.7\) & 13.3 & 14.2 \\
\textbf{M4} & \textbf{\(+\,L_{\mathrm{pop}}\)} & \textbf{8} & \textbf{0.998} & \textbf{\(-155.6\)} & \textbf{\(-148.0\)} & \textbf{\(-144.9\)} & \textbf{0.0} & \textbf{0.0} \\
\bottomrule
\end{tabular}
\end{table}

% ================================================================
\section{Forecast}

The forecast window
\([2028.4,2034.4]\)~CE
should be read as the candidate interval in which the present
\(H{=}6\)
regime must reveal whether it is genuinely non-terminal.
At the level of the present ontology,
there are only three broad outcomes:
terminal failure,
stationary arrest,
or successful advance.
The 2031.4~CE central date is therefore not an isolated prediction.
It is the point at which those possibilities become empirically separable.

To avoid confusion with the historical Col1/Col2 collapse labels already used for Roman and medieval resets,
we denote the prospective alternatives as route I,
route II,
and route III.

\begin{table}[htbp]
\centering
\caption{
Prospective alternatives for the present \(H{=}6\) frontier.
These are future-route labels,
not the historical Col1/Col2 collapse windows already realized within
\(\Omega_{\mathrm{social}}\).
}
\label{tab:forecast_routes}
\begin{tabular}{p{2.3cm}p{4.1cm}p{7.0cm}}
\toprule
Route & Structural reading & Meaning in the ladder ontology \\
\midrule
I & Malthusian collapse & demographic/resource exhaustion blocks further relay advance and forces failure before a new rung stabilizes \\
II & Technological arrest & \(H{=}6\) behaves as a stable endpoint, so the compression law terminates and Eq.~\eqref{eq:h_accel} is falsified \\
III & Ladder inversion / retrogression & the present carrier loses relay capacity and the system falls back toward \(H\le 5\) rather than sustaining the current frontier \\
\bottomrule
\end{tabular}
\end{table}

These three routes are exhaustive at the reduced-form level used here.
If route~I is realized,
the ladder fails by collapse.
If route~II is realized,
the ladder fails by stable arrest.
If route~III is realized,
the ladder fails by inversion.
Only when all three are excluded does the argument become one of structural necessity:
a further rung is then the only surviving continuation of the historical ontology.
That inference is strong,
but it is not substrate-specific.
It says that \(H{=}6\) cannot be the last rung if the system neither collapses,
freezes,
nor retreats.
It does not say what material or institutional form the next rung must take.
The prospective route logic and its reject-if-observed clauses are deferred to SI Appendix~\ref{sec:S-Collapse3Routes}.

The present paper therefore makes a sharply falsifiable forecast.
Persistence of the current rung beyond the upper bound of the window,
without discrete advance,
counts against the secular acceleration law.
By contrast,
an observable transition inside the window would support the claim that
\(H{=}6\)
was a metastable frontier rather than a civilizational endpoint.
The paper is thus not committed to optimism.
It is committed to a narrow structural claim:
late-history timing is precise enough that the ontology can now be tested prospectively.

% ================================================================
\section{Discussion}

The central gain of the V64 reframe is that the paper no longer reads as a loose mixture of singularity theory,
demographic transition,
and speculative frontier timing.
It reads instead as a single structural argument.
The 70,000-year record is organized by one pairwise backbone,
its realized states are discrete,
and the discreteness is not imposed after the fact.
The pre-H3 stall corridor,
the exogenous-ceiling theorem,
the two realized historical collapse windows,
and the milestone cross-check all converge on the same conclusion:
civilization history is better read as ladder ontology than as a smooth one-parameter growth curve.

This also clarifies the paper's place among existing traditions.
Malthusian thinking is recovered whenever the active constraint is food capacity.
Unified growth theory is recovered whenever population pressure forces a ceiling replacement.
Cliodynamic collapse theory is recovered whenever social-friction overshoot dominates.
What the present paper adds is the state variable those traditions usually leave implicit:
the rung itself.
Once rung selection is made explicit,
"why the singularity failed" and
"why the next ceiling matters"
become the same question.

The V64 framing is intentionally narrower than V63 on one point.
It does not attempt to identify the next carrier in the main text.
That is not because the question is unimportant.
It is because the present paper can already make a strong claim without answering it.
The non-terminal status of
\(H{=}6\)
follows from the joint structure of the ladder,
the waiting-time compression law,
and the collapse/arrest alternatives.
That is the load-bearing result.
Everything else is a realization problem layered on top of it.

\subsection*{Open question: the \(H{=}7\) substrate}\label{sec:open-h7-substrate}

The present paper therefore stops at a deliberate boundary.
It argues that a prospective
\(H{=}7\)
regime is structurally required if the ladder continues,
but it does not identify the substrate that would support that rung.
The candidate space is intentionally left broad:
biological uplift,
a technological substrate,
a hybrid carrier,
or some other presently unclassified realization all remain admissible at the level of this paper.
The argument here is only that the ontology has another rung,
not that the realizing material is already known.
Detailed substrate analysis is the subject of Paper~5~\cite{Kato2026e}.

\subsection*{Falsifiability}

The reframe sharpens rather than weakens falsifiability.
The ladder fails if the 70,000-year record is better described by a smooth closure with no discrete gate structure,
if the exogenous ceiling can be endogenized without knife-edge instability and collapse-count explosion,
if the Humanness theorem is broken by a lineage lacking one of
\(C_1\),
\(C_2\),
and
\(C_3\),
or if the present
\(H{=}6\)
state persists well beyond the forecast window without either discrete advance or a route-I/II/III failure pattern.
That is a demanding standard,
but it is also what makes the present ontology more than a metaphor.
It is a historical structure with explicit reject-if-observed clauses.

% ================================================================
\section{Methods}

\paragraph{Data.}
The 28-point calibration panel is a harmonized anchor-point reconstruction rather than a uniform census series.
Its historical scaffold is the corrected Biraben world-population table,
cross-checked against McEvedy and Jones for consistency of premodern levels and turning points;
late modern anchors are taken from United Nations World Population Prospects 2022~\cite{biraben1979,mcevedy1978,unwpp2022}.
HYDE 3.2 is used only as an external robustness comparator for 10,000~BCE--1850~CE,
not to construct the fitted panel,
and the Thomlinson estimate at
\(-\)25,000~BCE is used only for out-of-sample validation~\cite{goldewijk2017,thomlinson1975}.
The documented cleaning step is downward revision of three early Biraben anchors
(\(-\)8000:
8\(\to\)5~M;
\(-\)5000:
15\(\to\)5~M;
\(-\)3000:
30\(\to\)14~M)
to maintain consistency with the hunter-gatherer ceiling
\(K_{\mathrm{HG}}=5\)~M,
with all values expressed in millions;
the full correction log with provenance is archived in the reproducibility repository.
The model is fit to the discrete 28 anchors themselves,
not to an interpolated annual series.
Premodern anchors should therefore be read as reconstructed estimates rather than direct observations:
pre-agricultural points
(\(-\)70k to \(-\)10k)
carry factor 2--5 uncertainty and are used qualitatively only~\cite{hassan1981,ambrose1998},
whereas retained pre-1500 anchors are treated conservatively as roughly factor-of-two uncertain;
post-1800 UN-based anchors are substantially tighter.

\paragraph{Model.}
The backbone law is Eq.~\eqref{eq:backbone},
\(dN/dt=a_0N^2\),
with
\(a_0=4.74\times 10^{-12}\)
and
\(N\)
measured in millions.
The model is organized in two levels.
Level~1,
which is the level analyzed in the present paper,
uses the exogenous reduced-form social-friction envelope
\begin{equation}
V_F(t)=C(t_c-t)^{-\alpha},
\end{equation}
with fitted
\((C,\alpha,t_c)\)
and threshold dynamics governed by
\((x_{\mathrm{on}},x_{\mathrm{off}})\).
This exogeneity claim is falsifiable rather than definitional:
any endogenous closure
\(V_F=f(N,\Omega)\)
that simultaneously
(a) reproduces exactly two collapse episodes,
(b) matches the 28-point calibration at
\(R^2>0.97\),
and
(c) retains Jacobian condition number
\(\kappa(J)<10^4\)
would refute the theorem-plus-simulation package reported here.

To represent the post-1963 fertility transition in reduced form,
the full model introduces a time-dependent demographic multiplier
\(L_{\mathrm{pop}}(t)\).
We set
\(L_{\mathrm{pop}}(t)=1\)
for
\(t<1963\).
For
\(t\ge 1963\),
\(L_{\mathrm{pop}}\)
follows a logistic decline toward a long-run floor,
\begin{equation}
\frac{dL_{\mathrm{pop}}}{dt}
=r_L\,L_{\mathrm{pop}}(t)\!\left(1-\frac{L_{\mathrm{pop}}(t)}{L_{\mathrm{floor}}}\right),
\qquad
L_{\mathrm{pop}}(1963)=1,
\label{eq:lpop_ode}
\end{equation}
with
\(L_{\mathrm{floor}}=0.52\)
and
\(r_L=0.01328\),
implying a 75-year characteristic transition timescale consistent with post-1963 global data.
When
\(L_{\mathrm{pop}}<1\),
the effective exponent of the backbone law drops to
\begin{equation}
p_{\mathrm{eff}}(t):=p\times L_{\mathrm{pop}}(t),
\label{eq:peff}
\end{equation}
so that
\(p_{\mathrm{eff}}\to 2\times 0.52=1.04\)
as
\(t\to\infty\).
This explains the apparent post-1950 drop from
\(p=2\)
to
\(p\approx 1\)
without any structural exponent change:
the demographic transition suppresses the multiplicative growth factor,
not the interaction law itself.
The collapse-enabled full model is
\begin{equation}
\frac{dN}{dt}
=a_0\,N(t)^2L_{\mathrm{pop}}(t)
\left(1-\frac{N(t)}{\min\{K(t),V_F(t)\}}\right)\chi(t),
\label{eq:full_model}
\end{equation}
where
\(\chi(t)\in\{0,1\}\)
is the collapse-state indicator governed by hysteretic thresholds.
The main text reports the simpler exogenous-\(V_F\) run with
\(L_{\mathrm{pop}}=1\)
throughout;
the
\(L_{\mathrm{pop}}\)-enabled specification is retained as the full model used for validation of the modern slowdown.

\paragraph{Validation.}
Validation uses six independent checks.
First,
segment-wise RMS decomposition prevents single-segment overfitting;
the preserved decomposition is HG:
0.000,
Food:
0.022,
Social:
0.056,
Modern:
0.034.
Second,
the model predicts
\(N_{\mathrm{pred}}=3.1\)~M at
\(-\)25,000~BCE versus
\(N_{\mathrm{obs}}=3.3\)~M,
an out-of-sample error below 7\%.
A systematic out-of-sample test with a pre-1800 training cutoff
(18 pre-industrial points)
yields
\(R^2_{\mathrm{test}}=0.81\)
and predicts 2024 world population within 3.1\%.
Third,
cross-dataset robustness is checked against HYDE 3.2 for 10,000~BCE--1850~CE~\cite{goldewijk2017}.
Fourth,
free-\(p\)
estimation over 1000~BCE--1963~CE yields
\(p=2.14\),
while
\(\Delta\mathrm{AIC}=+383\)
versus exponential growth provides decisive evidence for the
\(N^2\)
backbone.
Fifth,
leave-one-out cross-validation yields median absolute percentage error 3.5\% and mean 7.4\%,
confirming that no single calibration point drives the fit.
Sixth,
nested AIC\(_c\)/BIC model selection confirms that the full model
(M4,
\(k=8\) parameters)
is optimal,
with
\(\Delta\mathrm{AIC}_c=13.3\)
relative to M3.

\paragraph{Numerics.}
Collapse begins when
\(N/V_F\ge x_{\mathrm{on}}\)
and ends when
\(N/V_F\le x_{\mathrm{off}}\).
During collapse,
the main reported run freezes population at
\(dN=0\).
The endogenous-\(V_F\) falsification family used in the SI is sampled from smooth positive monotone closures and accepted only when solver convergence and collapse counts are stable under step-halving and tighter tolerances.

\paragraph{Statistics.}
All reported comparisons include effect sizes
(\(R^2\) for variance explained,
\(\Delta\mathrm{AIC}_c\) and
\(\Delta\mathrm{BIC}\) for nested model selection),
exact sample sizes
(\(n=28\) calibration anchors unless stated otherwise),
95\% confidence intervals from case-bootstrap resamples,
and exact
\(P\)
values from sequential
\(F\)-tests.
Cohen's
\(d\)
equivalents are not applicable to single-trajectory reduced-form fits;
the primary effect-size metric is the proportional variance explained and the model-comparison information loss.

% ================================================================
\section*{Acknowledgments}
The author thanks the anonymous reviewers for their comments.

\bibliographystyle{unsrt}
\bibliography{paper3_refs_V13}

\section*{Author contributions}
S.K. conceptualized the study,
developed the methodology,
performed formal analysis,
wrote the original draft,
and prepared all visualizations.

\section*{Competing interests}
The author declares no competing interests.

\section*{Data availability}
Historical population reconstructions used in this study are drawn from Biraben (1979, 2003),
McEvedy and Jones (1978),
and the United Nations World Population Prospects 2022.
The corrected 28-point calibration table with full correction provenance,
the Python simulation code,
figure-generation notebooks,
and all parameter files will be deposited in a Zenodo repository upon acceptance
(DOI:
\texttt{10.5281/zenodo.XXXXXXX}).
No proprietary data were used.

\section*{Code availability}
The Python simulation code,
figure-generation notebooks,
and shell scripts to reproduce all main-text and SI figures will be deposited in the same Zenodo repository
(DOI:
\texttt{10.5281/zenodo.XXXXXXX}).

\section*{Figure legends}

\paragraph{Figure~3 $\vert$ Species ladder and the \(H{=}3\) threshold.}
Figure~3 visualizes the Humanness theorem as a species-level classification problem.
The key point is not the exact continuous coordinate of any one non-human lineage,
but the hard threshold at
\(H{=}3\):
without the full conjunction
\(C_1\wedge C_2\wedge C_3\),
no lineage reaches the civilizational relay axis.

\paragraph{Figure~4 $\vert$ Historical timeline and the non-terminal status of the present rung.}
Figure~4 aligns the six realized
\(\Omega\)-regimes with the relay-ceiling reading of the historical record.
The dashed future extension should be read as a prospective rung boundary,
not as a substrate-specific claim.
Its role is to make visible the non-terminal status of the present
\(H{=}6\)
regime under Eq.~\eqref{eq:h_accel}.

\section*{Pre-registration}
The present study is an analysis of historical reconstructed population data;
no prospective experimental protocol was registered,
and pre-registration is therefore not applicable.
All analysis choices are stated explicitly in the main text and SI and are reproducible from the deposited code.

\end{document}
% （以上で Paper 3 main V64 Kato Ladder Civilizational Ontology reframe 本文 終了）
